\section{Background and Motivation}
\label{sec:background}

\subsection{GUI effects and mobile interruptions}
\label{sec:bg-effects}
A GUI agent observes the screen through screenshots and the accessibility tree, and acts by injecting taps, text and gestures,
typically in a ReAct-style loop driven by an LLM or VLM~\cite{yao2023react,rawles2025androidworld,qin2025uitars}. Most actions only navigate.
A few cross a \emph{commit point}: after ``Send'', ``Pay'', ``Save'' or ``Post'' the world has changed whether or not the agent sees the
confirmation. We call the action sequence that ends at a commit point an \emph{effect}, and target effects whose only interface is the GUI.

Research agents (M3A, T3A, UI-TARS, Mobile-Agent) are host processes that drive a phone over adb~\cite{rawles2025androidworld}. On-device agents are apps with an
accessibility service~\cite{wen2024autodroid}. Interruptions are of four kinds. (1)~Agent death: a host agent crashes, or the low-memory killer reaps an on-device
agent~\cite{android_lmkd}. (2)~App death: the low-memory killer reaps the foreground app last~\cite{android_lmkd}, so the realistic case is \emph{background-then-kill}. Even that may be
rare for an app that a sync service keeps perceptible: \texttt{am kill} spared Element in every trial (\S\ref{sec:eval-real}), so for such apps agent death is the likelier
interruption. Background-then-kill keeps scheduled
WorkManager jobs; a crash or force-stop cancels alarms and jobs~\cite{android_workmanager}, and unpreserved state is lost~\cite{farooq2020livedroid}. (3)~Deferral:
App Standby, background restrictions and Doze batch network access and defer jobs~\cite{android_doze,android_standby}. (4)~Connectivity loss~\cite{sentosa2025cellreplay}:
an app may send a request and lose the acknowledgement. Each can leave the effect absent, committed, or still in flight.

\subsection{Why API-level recovery does not transfer}
\label{sec:bg-gap}
\begin{table}[!htbp]
\centering
\caption{Recent recovery work. Key: idempotency key (opt.: with and without). Read-back: D exact, EC eventually consistent, set partial, N noisy. Ask: appr.\ approval gate, part.\ exposes ambiguity. Crash: survives agent death. $^c$: concurrent work.}
\label{tab:position}
\footnotesize
\setlength{\tabcolsep}{1.7pt}
\resizebox{\columnwidth}{!}{%
\begin{tabular}{@{}l ccccc l@{}}
\toprule
Work & Key & Read & Ask & Crash & GUI & Guarantee / output \\
\midrule
Li / LIMBO$^c$~\cite{li2026exactlyonce} & opt. & D/EC & no & no & no & impossibility, measurement \\
RetryBench$^c$~\cite{liu2026retrybench} & opt. & D & no & no & no & retry-permission scoring \\
UndoBench$^c$~\cite{sah2026undobench} & opt. & D & no & no & no & recovery benchmark \\
Evidence-guided$^c$~\cite{zhu2026evidence} & opt. & set & yes & no & no & completion certificate \\
Atomix~\cite{mohammadi2026atomix} & -- & D & no & partial & no & transactions \\
Agent libOS~\cite{zhang2026agentlibos} & -- & D & part. & yes & no & no blind replay \\
Cordon~\cite{chen2026cordon} & -- & D & appr. & yes & no & staged effects \\
Resume Contract~\cite{khan2026resume} & -- & -- & no & yes & no & TLA+ contract \\
ACRFence~\cite{zheng2026acrfence} & -- & -- & no & yes & no & replay-or-fork \\
VeriGUI~\cite{zhang2026verigui} & no & N & no & no & yes & learned verifier \\
TaskVM~\cite{yang2026taskvm} & no & N & yes & -- & yes & task-state contract \\
AgentProg~\cite{tian2025agentprog} & no & -- & no & no & yes & redo \\
\textbf{\sys} & \textbf{no} & \textbf{N} & \textbf{yes} & \textbf{yes} & \textbf{yes} & $(\alpha,\beta)$; sim.\ + emulator \\
\bottomrule
\end{tabular}}
\end{table}

Table~\ref{tab:position} summarizes recent recovery work. Three gaps matter.

The first gap is the missing key. Idempotency keys~\cite{stripe_idempotency,helland2012idempotence} let a server collapse retries, and agents use them when they exist~\cite{li2026exactlyonce}.
A GUI ``Send'' button takes no key from the agent; the app may attach its own request id (\S\ref{sec:twohop}), but a GUI redo creates a new request with a new id.

The second gap is perceptual read-back. LIMBO's read paths are exact, eventually consistent, or missing~\cite{li2026exactlyonce}. Verified tool calls~\cite{mansoor2026verified},
RetryBench~\cite{liu2026retrybench}, explicit effect semantics~\cite{wang2026callability} and evidence-guided recovery~\cite{zhu2026evidence} reason over exact, possibly partial
evidence and keep unknown outcomes explicit. A GUI read-back is a perception that can be wrong, tends to repeat its error on a re-read, and is defeated by look-alikes;
VeriGUI trains an action-effect verifier for this setting~\cite{zhang2026verigui} without an error guarantee.

The third gap is the death of the agent and late commits. Li injects faults at the service boundary~\cite{li2026exactlyonce}, but on a phone the agent itself may die. Atomix keeps progress
in its runtime~\cite{mohammadi2026atomix}, and LangGraph is at-least-once across crashes~\cite{khan2026resume}. Agent libOS and Cordon stage effects behind durable intent
(prepare, dispatch, settle)~\cite{zhang2026agentlibos,chen2026cordon}. Our ledger follows that discipline, but settlement must be decided from perception.

\subsection{How often do agents create ghost effects?}
\label{sec:measure-teaser}
No published study measures duplicate side effects of mobile GUI agents under interruption. D-GARA evaluates Android agents under anomalies such as permission dialogs and
update prompts and reports degraded task success~\cite{chen2025dgara}; for tool agents, Li reports 56 to 74\% duplicates when a read-back cannot reveal the outcome~\cite{li2026exactlyonce},
and UndoBench 53.33\% for naive retry~\cite{sah2026undobench}. Our emulator measurements (\S\ref{sec:eval-real}) show the trap and its outcome: after a lost reply K-9~Mail showed every
accepted mail as unsent, and a scripted redo of the step (recompose, send) duplicated \RealGuiDup/\RealGuiN{} mail and calendar effects. This is the client-side face of the
SMTP duplicate problem known since 1988~\cite{rfc1047}. An LLM agent takes that step: \AgentFramework{} on \AgentModel{} resent the mail in \AgentResentN{} of \AgentValidN{} valid
episodes, each time leaving a duplicate at the server (LLM-agent experiment, \S\ref{sec:eval-agent}). T3A, a second agent, resent in
\AgentTextResentN{} of \AgentTextValidN. How often other agents and apps do so remains future work (\S\ref{sec:device-plan}).
